
     \documentclass[fleqn]{article}      
      \usepackage{amsmath}
      \usepackage{fontspec}
      \usepackage{kotex} % 한국어 지원  

      \begin{document}      
       

쉬움: \\\\ 
1. b) 함수는 입력에 대해 출력이 정해진 규칙입니다. \\\\
2. a) $y = 2x + 3$은 함수의 형태를 가지고 있습니다. \\\\
3. c) 3 - $f(2) = 2 + 1 = 3$입니다. \\\\
  
중간: \\\\  
1. c) 둘 다 - 함수의 그래프는 직선일 수도 있고 곡선일 수도 있습니다. \\\\
2. b) 포물선 - $f(x) = x^2$의 그래프는 포물선의 형태입니다. \\\\
3. b) 하나의 입력이 여러 개의 출력을 가질 수 없음 - 함수는 각 입력에 대해 오직 하나의 출력이 있어야 합니다. \\\\

어려움: \\\\  
1. a) 입력과 출력을 바꾼 함수 - 역함수는 원래 함수의 입력과 출력을 서로 바꾼 것입니다. \\\\
2. a) $f^{-1}(x) = \frac{x-2}{3}$ - 역함수를 구하기 위해서는 $y = 3x + 2$를 $x$에 대해 풀어야 합니다. \\\\
3. a) 3 - $x^3 = 27$일 때, $x = 3$입니다. \\\\      
      \end{document} 
  